% ============================================================
\chapter{CVaR で最適化する}\label{ch:cvar}

BL の事後の期待リターンができたら、比率を決める最適化が必要である。ふつうは分散をリスクとするが、裾の損失を気にするなら、CVaR（条件付きバリュー・アット・リスク）をリスクにできる。調査メモも、景気の局面のビューと CVaR を組み合わせる研究や、モメンタムの急落に CVaR で備える方法に触れている。この章では、CVaR の最適化の仕組みと、BL と組み合わせるときの注意をまとめる。

\section{VaR と CVaR}

ポートフォリオの損失を $L$ とする。信頼水準 $\beta$（例えば 95\%）の VaR は、損失が $1-\beta$ の確率でしか超えない水準である。CVaR は、損失が VaR を超えたときの損失の平均である。
\begin{equation}
\VaR_\beta=\min\{\ell:\Pr(L\le\ell)\ge\beta\},\qquad
\CVaR_\beta=\E[L\mid L\ge\VaR_\beta]
\end{equation}
この式は分布が連続な場合のもので、シナリオのような離散の分布では、「損失の大きい側の確率 $1-\beta$ の部分の平均」と定義する（次の節の式はこの定義と一致する）。VaR は「どこまで損するか」の境目しか見ないが、CVaR は境目を超えた先の大きさも見る。さらに、VaR は 2 つのポートフォリオを合わせたときのリスクが、それぞれのリスクの和を超えることがあり（劣加法性を満たさない）、分散投資の効果を正しく表さない。CVaR は劣加法性を含む「整合的なリスク尺度」の条件（Artzner ほか 1999）を満たす。

\section{Rockafellar と Uryasev の定式化}

Rockafellar と Uryasev（2000）は、CVaR を次の最小化の値として書けることを示した。
\begin{equation}
\CVaR_\beta(w)=\min_\zeta\Bigl\{\zeta+\frac{1}{1-\beta}\E\bigl[(L(w)-\zeta)^+\bigr]\Bigr\}
\label{eq:ru}
\end{equation}
最小にする $\zeta$ は VaR である。リターンのシナリオ $x_1,\dots,x_J$ と確率 $q_j$ があれば、損失は $L_j=-x_j\T w$ で、CVaR の最小化は線形計画になる。
\begin{equation}
\min_{w,\zeta,u}\;\zeta+\frac{1}{1-\beta}\sum_j q_ju_j
\quad\text{s.t.}\quad u_j\ge-x_j\T w-\zeta,\quad u_j\ge0,\quad \sum_iw_i=1,\quad w\ge0
\label{eq:ru-lp}
\end{equation}
期待リターンの下限や比率の上限も、線形の制約として足せる。変数はシナリオの数だけ増えるが、線形計画なので数万のシナリオでも解ける。

\section{BL と組み合わせる}

BL と CVaR の組み合わせ方は 2 つある。

\begin{enumerate}
\item \textbf{期待リターンだけ BL}：BL の事後の期待リターン $\muBL$ を目的関数や制約に使い、CVaR はリスクのモデル（過去のリターン、$t$ 分布、GARCH とコピュラのシミュレーション）のシナリオで測る。Teplova ほか（2023）、Alpern と Rachev（2025）はこの形である。Alpern と Rachev は、$\alpha\,w\T\muBL-(1-\alpha)\CVaR$ に回転率の罰則を足した目的関数を使っている。
\item \textbf{事後の分布ごと}：エントロピー・プーリングや COP で、ビューを入れた事後のシナリオ（確率）を作り、そのシナリオで期待リターンも CVaR も測る。
\end{enumerate}

\section{数値例}

\chapref{ch:ep} のエントロピー・プーリングの事後の確率（3 つのビューを入れたもの）のもとで、期待リターンをベンチマーク（セクターの時価総額の比率）と同じ月 \V{cv.target}\% 以上に保ちながら、分散を最小にする比率と、CVaR（95\%）を最小にする比率を比べた。比率の上限は、Teplova ほかにならって 30\%、40\%、50\% を試した。

\begin{table}[htbp]
\centering
\caption{エントロピー・プーリングの事後の確率のもとでの、分散の最小化と CVaR の最小化（月次 \%）。期待リターンはベンチマークと同じ水準以上。下の 4 行は事後の確率で測ったリスク。}
\label{tab:cvar}
\tabin{cvar}
\end{table}

\tabref{tab:cvar} から分かることを挙げる。

\begin{itemize}
\item 分散の最小化（上限 40\%）と CVaR の最小化（上限 40\%）の比率は似ている。CVaR の最小化は、CVaR を \V{cv.mv.cvar}\% から \V{cv.cv.cvar}\% にわずかに下げ、標準偏差はほとんど変わらなかった（\V{cv.mv.sd}\% と \V{cv.cv.sd}\%）。上限 50\% は効いておらず、40\% と同じ比率になった。
\item どちらも、ベンチマークより CVaR と標準偏差がかなり小さい。期待リターンを同じに保ちながら、リスクの小さいセクター（生活必需品など）と、期待リターンの高いセクター（情報技術）の組み合わせに寄せている。
\item 上限を 30\% に下げると、情報技術が上限に張り付き、残りが別のセクターに分かれる。上限は比率の形を大きく変える。
\end{itemize}

\section{楕円分布では、分散と CVaR は同じ答えになる}

分散の最小化と CVaR の最小化が似た比率になったのは偶然ではない。リターンが正規分布や $t$ 分布のような楕円分布に従うなら、ポートフォリオの CVaR は
\begin{equation}
\CVaR_\beta(w)=-w\T\mu+k_\beta\sqrt{w\T\Sigma w}
\end{equation}
の形になる（$k_\beta$ は分布と $\beta$ だけで決まる定数）。期待リターンを固定すれば、CVaR の最小化は標準偏差の最小化と同じになる。今回の事前分布は $t$ 分布で、エントロピー・プーリングで確率を少し変えたため、事後は楕円分布からわずかにずれ、答えもわずかにずれた。

CVaR を使う意味があるのは、分布が非対称なとき（オプション、信用リスク、モメンタムの急落）や、裾の依存が正規分布と違うときである。株式のセクターのように、おおむね楕円分布とみなせる資産では、分散の最小化で十分なことが多い。

\section{推定の誤差}

CVaR は、分布の裾の少数のシナリオで決まる。95\% の CVaR なら、全体の 5\% のシナリオの平均である。推定に使えるシナリオが少ないと、CVaR の推定は不安定で、最適な比率も大きく揺れる。Alpern と Rachev は、BL の期待リターンと CVaR を組み合わせた戦略で回転率が大きくなりすぎ、回転率の制約を強めて対処している。Teplova ほかは、コピュラ-CVaR の戦略のほうが回転率が低く、取引費用を考えた損益分岐点で有利だったと報告している。

\section{調査メモの CVaR-BL-MCV について}

調査メモは、景気循環のビュー（MCV）と CVaR を組み合わせた「CVaR-BL-MCV」が、年金の資産配分でシャープレシオを平均 19.7\% 改善した、と紹介し、CVaR の定義式を示している。定義式は Rockafellar と Uryasev の式~\eqref{eq:ru} と同じ形である。モデルと数字の出典は、この資料では確認できなかった（\appref{app:check}）。

\begin{tipbox}[コツ：CVaR を使うとき]
\begin{enumerate}
\item 期待リターンは BL で作り、CVaR はリスクのモデルのシナリオで測る、と分けて考える。
\item 楕円分布に近い資産では、分散の最小化とほとんど同じ答えになる。非対称な資産や、裾の連動が強い局面で使う。
\item 裾のシナリオの数を十分に取る（95\% の CVaR なら、裾に数百以上）。
\item 比率が揺れやすいので、回転率の制約や罰則を一緒に入れる。
\end{enumerate}
\end{tipbox}

\begin{keybox}[この章のまとめ]
\begin{itemize}
\item CVaR は、VaR を超えた損失の平均で、整合的なリスク尺度である。Rockafellar と Uryasev の定式化で、シナリオから線形計画として最小化できる。
\item BL と組み合わせるときは、期待リターンを BL で、リスクを裾の厚いシナリオで測るのが基本である。
\item 楕円分布なら、CVaR の最小化は分散の最小化と同じ答えになる。例でも、両者の比率は近かった。
\end{itemize}
\end{keybox}
